\section{Purpose and Method}

\subsection{Why an algebra}

\texttt{fsb}'s central promise is negative: a path that the user has denied is
never listed, searched, previewed, downloaded, or even confirmed to exist.  That
promise is a statement about \emph{every} way of reaching a file, so it is
threatened whenever two code paths answer the same question differently.  The
defects found and fixed in the project so far had exactly that shape: a listing
through a symlinked directory that leaked denied names, a trailing newline that
defeated a rule, an unreadable file answering 403 instead of 404, and a macOS
firmlink alias that bypassed a rule.

An algebraic specification forces the question to be written down once, and
then every operation to be compared with it.  This document does that for four
parts of the system:

\begin{enumerate}
  \item the rule matcher (Section~\ref{sec:rules});
  \item the access decision shared by all endpoints (Section~\ref{sec:access});
  \item the content-typed endpoints and the pure frontend functions
        (Sections~\ref{sec:content} and~\ref{sec:frontend}); and
  \item the preview state machine (Section~\ref{sec:state}).
\end{enumerate}

\subsection{Method}

Each law is stated in the form
\emph{for all inputs of a stated kind, a stated equation or implication holds},
then given an executable check: a randomized property test with a fixed seed, a
test against a fixture that assembles the awkward cases together, or an
end-to-end browser test.  Laws that failed when first checked are recorded in
Section~\ref{sec:findings} together with the fix.  Laws that hold only under a
stated assumption say so, and the assumptions that cannot be discharged inside
the program are collected as \emph{boundaries}.

\subsection{What is not covered}

The document does not model layout, streaming, timing, or the browser's own
behavior (except where the program relies on it).  It does not replace an
independent review; it is intended to give a reviewer a precise statement of
what the program claims.
